% =============================================================================
\section{Modelo microscópico de la polarización}\label{sec:microscopico}
% =============================================================================

Vamos a describir los dieléctricos mediante la \emph{aproximación dipolar}: un dieléctrico polarizado es,
a escala macroscópica, una distribución de dipolos. Antes de adoptarla conviene justificarla con un
modelo microscópico sencillo, que además nos dará los órdenes de magnitud. Hay que distinguir dos clases
de sustancias según que sus moléculas tengan o no momento dipolar en ausencia de campo.

\subsection{Sustancias no polares: el átomo de nube uniforme}\label{sec:no-polares}

En una sustancia no polar el momento dipolar de cada átomo o molécula es nulo, $\vp=\vcero$, a escala
atómica. El modelo más simple de átomo es un núcleo puntual de carga $+q$ rodeado de una nube electrónica
esférica de carga $-q$ uniformemente repartida en un radio $a$ (\cref{fig:atomo}). Sin campo aplicado,
el centro de la nube coincide con el núcleo y el momento dipolar es nulo.

% Figura: modelo de nube electrónica uniforme (átomo sin polarizar / polarizado)
% y campo de la nube en función de la distancia a su centro.
\begin{figure}[t]
\centering
\begin{subfigure}[b]{0.3\textwidth}\centering
\begin{tikzpicture}[scale=1.0]
  \shade[inner color=acentoA!38!fondo, outer color=acentoA!14!fondo] (0,0) circle (1.15);
  \draw[acentoA!80!fondo, line width=0.6pt] (0,0) circle (1.15);
  \draw[tenue, dashed, line width=0.4pt] (0,0) -- (-40:1.15) node[midway, below left=-1pt, etiq] {$a$};
  \fill[cargapos] (0,0) circle (0.1);
  \node[etiq, above=2pt] at (0,0.05) {$+q$};
  \node[etiq] at (0.72,0.72) {$-q$};
  \node[etiq] at (0,1.55) {$\vE=\vcero$};
\end{tikzpicture}
\caption{Átomo sin polarizar}
\end{subfigure}\hfill
\begin{subfigure}[b]{0.3\textwidth}\centering
\begin{tikzpicture}[scale=1.0]
  \shade[inner color=acentoA!38!fondo, outer color=acentoA!14!fondo] (0,0) circle (1.15);
  \draw[acentoA!80!fondo, line width=0.6pt] (0,0) circle (1.15);
  \fill[texto] (0,0) circle (0.035);
  \draw[campoP, line width=0.9pt] (0,0) -- (0.42,0);
  \node[etiq, below=1pt] at (0.21,0) {$d$};
  \fill[cargapos] (0.42,0) circle (0.1);
  \node[etiq, above=2pt] at (0.42,0.05) {$+q$};
  \node[etiq] at (-0.62,0.72) {$-q$};
  \draw[vecE] (-0.7,1.55) -- (0.7,1.55) node[right, etiq] {$\vE$};
\end{tikzpicture}
\caption{Átomo polarizado}
\end{subfigure}\hfill
\begin{subfigure}[b]{0.36\textwidth}\centering
\begin{tikzpicture}
\begin{axis}[width=5.6cm, height=4.1cm, axis lines=left, xmin=0, xmax=3.2, ymin=0, ymax=1.18,
  xlabel={$r/a$}, ylabel={$E_{\text{nube}}\big/\tfrac{q}{4\pi\eo a^{2}}$},
  xtick={0,1,2,3}, ytick={0,0.5,1}, ylabel style={yshift=-3pt}]
  \addplot[curvaE, domain=0:1, samples=2] {x};
  \addplot[curvaE, domain=1:3.2, samples=80] {1/x^2};
  \draw[equipot] (axis cs:1,0) -- (axis cs:1,1.18);
  \node[etiqt, anchor=west] at (axis cs:0.08,0.86) {$\propto r$};
  \node[etiqt, anchor=west] at (axis cs:1.75,0.52) {$\propto 1/r^{2}$};
\end{axis}
\end{tikzpicture}
\caption{Campo de la nube}
\end{subfigure}
\caption{Modelo atómico de nube uniforme. (a) Sin campo aplicado, el centro de la nube electrónica
(carga $-q$, radio $a$) coincide con el núcleo y el momento dipolar es nulo. (b) Un campo externo
desplaza el núcleo una distancia $d\ll a$ respecto al centro de la nube, y aparece $\vp=q\vec d$.
(c) Campo de la propia nube: dentro crece linealmente con $r$ (ley de Hooke electrostática),
fuera decae como el de una carga puntual; el equilibrio se alcanza en la rama lineal.}
\label{fig:atomo}
\end{figure}


Si aplicamos un campo externo $\vE$, el núcleo es empujado en el sentido del campo y la nube en el
contrario. Supondremos que la nube se desplaza rígidamente, sin deformarse, y que el desplazamiento
relativo $d$ es pequeño frente a $a$.

\begin{proposicion}[Polarizabilidad del átomo de nube uniforme]\label{prop:polarizabilidad}
En el modelo de nube uniforme, un campo externo $\vE$ induce el momento dipolar
\begin{equation}\label{eq:p-atomo}
  \vp=q\,\vec d=4\pi\eo a^3\,\vE ,
\end{equation}
donde $\vec d$ es el desplazamiento del núcleo respecto al centro de la nube.
\end{proposicion}

\begin{proof}
El campo de la nube en un punto interior a distancia $r$ de su centro se obtiene con la ley de Gauss: una
esfera concéntrica de radio $r<a$ encierra la fracción $(r/a)^3$ de la carga $-q$, luego
\[
  E_{\text{nube}}(r)\,4\pi r^2=\frac{-q}{\eo}\frac{r^3}{a^3}
  \quad\Longrightarrow\quad
  \vE_{\text{nube}}(\vr)=-\frac{q}{4\pi\eo a^3}\,\vr .
\]
Es un campo lineal en $\vr$ y dirigido hacia el centro (\cref{fig:atomo}\,(c)). Sobre el núcleo, situado en
$\vec d$, actúan el campo externo y el de la nube; el equilibrio exige
\[
  q\vE+q\,\vE_{\text{nube}}(\vec d)=\vcero
  \quad\Longrightarrow\quad
  \vE=\frac{q}{4\pi\eo a^3}\,\vec d ,
\]
y multiplicando por $4\pi\eo a^3$ se obtiene \eqref{eq:p-atomo}. La fuerza de la nube es recuperadora y
proporcional al desplazamiento, como la de un muelle: es la «ley de Hooke» del átomo.
\end{proof}

El resultado tiene la forma general de la respuesta lineal de un átomo isótropo.

\begin{definicion}[Polarizabilidad]\label{def:polarizabilidad}
Si el momento dipolar inducido en un átomo o molécula es proporcional al campo que actúa sobre él,
\begin{equation}\label{eq:alfa}
  \vp=\alpha\,\vE ,
\end{equation}
la constante $\alpha>0$ se llama \emph{polarizabilidad} (atómica o molecular). Sus unidades son
\[
  [\alpha]=\frac{\si{C\,m}}{\si{V/m}}=\frac{\si{C^2\,m}}{\si{N}}=\si{F\,m^2},
\]
pues $\si{C/V}=\si{F}$. En el modelo de nube uniforme, $\alpha=4\pi\eo a^3$: es proporcional al volumen atómico.
\end{definicion}

\begin{ejemplo}[El átomo de hidrógeno]\label{ej:hidrogeno}
Con el radio de Bohr, $a_0\simeq\SI{5.29e-11}{m}$, el modelo predice
\[
  \alpha_{\text{mod}}=4\pi\eo a_0^3\simeq 4\pi\,(\SI{8.854e-12}{F/m})(\SI{5.29e-11}{m})^3\simeq\SI{1.65e-41}{F\,m^2},
\]
mientras que el valor experimental es $\alpha_{\mathrm H}\simeq\SI{7.42e-41}{F\,m^2}$~\cite{Griffiths}: el modelo da
$\alpha_{\text{mod}}\simeq0{,}22\,\alpha_{\mathrm H}$. Para un modelo tan tosco, el orden de magnitud y la dependencia con
$a^3$ son correctos.
\end{ejemplo}

\begin{observacion}[El resultado cuántico]\label{obs:hidrogeno-cuantico}
La mecánica cuántica da para el estado fundamental del hidrógeno, de forma exacta,
$\alpha=\tfrac92\,4\pi\eo a_0^3$~\cite{Griffiths}. El factor $9/2$ es justamente el cociente $1/0{,}22$ del
\cref{ej:hidrogeno}: la nube real no es una esfera uniforme y rígida, sino que se deforma y es más
«blanda» que en nuestro modelo.
\end{observacion}

\begin{ejemplo}[Cuánto se desplaza la nube]\label{ej:desplazamiento}
Un campo de $E=\SI{5e5}{V/m}$ es muy intenso en el laboratorio, aunque inferior al de ruptura dieléctrica del
aire ($\sim\SI{3e6}{V/m}$~\cite{T3}). Sobre un átomo de hidrógeno induce
\[
  p=\alpha_{\mathrm H}E\simeq(\SI{7.42e-41}{F\,m^2})(\SI{5e5}{V/m})\simeq\SI{3.7e-35}{C\,m},
\]
y el núcleo se separa del centro de la nube una distancia
\[
  d=\frac{p}{e}\simeq\frac{\SI{3.7e-35}{C\,m}}{\SI{1.602e-19}{C}}\simeq\SI{2.3e-16}{m}=\SI{0.23}{fm}.
\]
El desplazamiento es del orden del tamaño de un núcleo y cinco órdenes de magnitud menor que el átomo: la
hipótesis $d\ll a$ está más que justificada, y con ella la respuesta lineal.
\end{ejemplo}

\subsection{Sustancias polares: orientación de dipolos permanentes}\label{sec:polares}

En una sustancia polar cada molécula tiene un momento dipolar permanente, $\vp\neq\vcero$, porque los
centros de carga positiva y negativa no coinciden. El ejemplo típico es el agua
(\cref{fig:agua}), con $p_{\mathrm{H_2O}}\simeq\SI{6.1e-30}{C\,m}$ \cite{Griffiths}, unas $10^5$ veces el
dipolo inducido del \cref{ej:desplazamiento}.

% Figura: la molécula de agua como dipolo permanente.
\begin{figure}[t]
\centering
\begin{subfigure}[b]{0.44\textwidth}\centering
\begin{tikzpicture}[scale=1.0]
  % geometría: ángulo H-O-H de 104,5 grados, dipolo según el bisector (eje +x)
  \coordinate (O) at (0,0);
  \coordinate (H1) at (52.25:1.25);
  \coordinate (H2) at (-52.25:1.25);
  \shade[ball color=carganeg!85!fondo] (H1) circle (0.52);
  \shade[ball color=carganeg!85!fondo] (H2) circle (0.52);
  \shade[ball color=cargapos!85!fondo] (O) circle (0.95);
  \node[font=\bfseries, text=white] at (-0.15,0) {$-2\delta$};
  \node[font=\bfseries, text=white] at ($(H1)+(0.12,0.08)$) {$+\delta$};
  \node[font=\bfseries, text=white] at ($(H2)+(0.12,-0.08)$) {$+\delta$};
  \draw[vecP, line width=1.3pt] (0.35,0) -- (2.15,0) node[right, etiq] {$\vp$};
  \node[etiqt] at ($(H1)+(0.95,0.25)$) {H};
  \node[etiqt] at ($(H2)+(0.95,-0.25)$) {H};
  \node[etiqt] at (-1.25,0.0) {O};
  \draw[tenue, line width=0.4pt] (52.25:0.55) arc (52.25:-52.25:0.55);
\end{tikzpicture}
\caption{Distribución de carga y dipolo}
\end{subfigure}\hfill
\begin{subfigure}[b]{0.52\textwidth}\centering
\begin{tikzpicture}[scale=0.95]
  \tikzset{electron/.style={circle, fill=acentoC, draw=fondo, line width=0.3pt, inner sep=1.5pt}}
  % oxígeno: núcleo y dos capas
  \draw[tenue, line width=0.5pt] (0,0) circle (0.62);
  \draw[tenue, line width=0.5pt] (0,0) circle (1.1);
  \shade[ball color=cargapos!85!fondo] (0,0) circle (0.36);
  \node[font=\tiny, text=white, align=center] at (0,0) {$8p$\\[-3pt]$8n$};
  \foreach \t in {90,270} \node[electron] at (\t:0.62) {};
  \foreach \t in {30,90,150,210,-30,-90} \node[electron] at (\t:1.1) {};
  % hidrógenos (enlaces covalentes: capas solapadas)
  \foreach \h/\ang in {A/-35,B/-145} {
    \coordinate (\h) at (\ang:1.82);
    \draw[tenue, line width=0.5pt] (\h) circle (0.62);
    \shade[ball color=carganeg!85!fondo] (\h) circle (0.2);
    \node[font=\tiny, text=white] at (\h) {$+$};
  }
  \node[electron] at ($(A)+(150:0.62)$) {};
  \node[electron] at ($(B)+(30:0.62)$) {};
  \node[etiqt] at (0,1.42) {oxígeno};
  \node[etiqt] at ($(A)+(0.1,-0.86)$) {H};
  \node[etiqt] at ($(B)+(-0.1,-0.86)$) {H};
\end{tikzpicture}
\caption{Estructura electrónica (esquema)}
\end{subfigure}
\caption{La molécula de agua, ejemplo de sustancia polar. (a) El oxígeno atrae con más fuerza los
electrones de los enlaces: queda una carga parcial $-2\delta$ en el oxígeno y $+\delta$ en cada hidrógeno.
Como la molécula no es lineal (ángulo de $104{,}5^\circ$), los centros de carga positiva y negativa no
coinciden y existe un momento dipolar permanente $p_{\mathrm{H_2O}}\simeq\SI{6.1e-30}{C\,m}$, dirigido
del oxígeno hacia los hidrógenos. (b) Esquema de capas: cada hidrógeno comparte su electrón con la
capa externa del oxígeno.}
\label{fig:agua}
\end{figure}


\begin{proposicion}[Dipolo en un campo uniforme]\label{prop:par}
Sobre un dipolo rígido $\vp=q\vec L$ en un campo uniforme $\vE$ la fuerza neta es nula y el momento de las
fuerzas respecto a cualquier punto es
\begin{equation}\label{eq:par}
  \vtau=\vp\times\vE ,
\end{equation}
que tiende a alinear $\vp$ con $\vE$. Su energía potencial es $U=-\vp\cdot\vE$, mínima cuando $\vp\parallel\vE$.
\end{proposicion}

\begin{proof}
Con las cargas $\pm q$ en $\vr_\pm$, la fuerza neta es $q\vE-q\vE=\vcero$ y el momento
$\vtau=\vr_+\times q\vE+\vr_-\times(-q\vE)=q(\vr_+-\vr_-)\times\vE=\vp\times\vE$, independiente del origen.
Si $\theta$ es el ángulo entre $\vp$ y $\vE$, $\abs{\vtau}=pE\sen\theta$ y el trabajo para girar el dipolo desde
$\theta=\pi/2$ es $\int_{\pi/2}^{\theta}pE\sen\theta'\,\dd\theta'=-pE\cos\theta=-\vp\cdot\vE$.
\end{proof}

Sin campo aplicado, la agitación térmica reparte las orientaciones al azar y el momento dipolar medio
de cualquier volumen macroscópico es nulo. Con campo, el par \eqref{eq:par} compite con la agitación
térmica y la orientación es solo parcial, pero el momento medio deja de anularse y tiene la dirección de
$\vE$ (\cref{fig:orientacion}).

% Figura: (a) par sobre un dipolo en un campo uniforme; (b) y (c) orientación
% térmica de dipolos permanentes sin campo y con campo.
\begin{figure}[t]
\centering
\begin{subfigure}[b]{0.38\textwidth}\centering
\begin{tikzpicture}[scale=1.0]
  \foreach \y in {-1.6,-1.07,...,1.62}
    \draw[campoE!55!fondo, line width=0.45pt, flechamedio=0.96] (-2.4,\y) -- (2.4,\y);
  \node[etiq, text=campoE] at (2.62,1.6) {$\vE$};
  \fill[dielec, draw=texto!60!fondo, line width=0.4pt, rotate=40] (0,0) ellipse (0.6 and 1.45);
  \coordinate (m) at (220:1.08); \coordinate (p) at (40:1.08);
  \draw[texto, line width=0.6pt] (m) -- (p);
  \draw[vecP, line width=1.2pt] (0,0) -- (40:0.7);
  \node[etiq] at (128:0.36) {$\vp$};
  \draw[texto!60!fondo, dashed, line width=0.35pt] (0,0) -- (1.05,0);
  \draw[tenue, line width=0.4pt] (0.5,0) arc (0:40:0.5);
  \node[etiqm, font=\footnotesize] at (17:0.72) {$\theta$};
  \fill[cargapos] (p) circle (0.085); \fill[carganeg] (m) circle (0.085);
  \draw[draw=cargapos, line width=1.0pt, -{Stealth[length=2.2mm]}] (p) -- ++(1.05,0);
  \node[etiqf, anchor=south] at ($(p)+(0.8,0.14)$) {$\vF_{+}=q\vE$};
  \draw[draw=carganeg, line width=1.0pt, -{Stealth[length=2.2mm]}] (m) -- ++(-1.05,0);
  \node[etiqf, anchor=north] at ($(m)+(-0.8,-0.14)$) {$\vF_{-}=-q\vE$};
  \node[etiqf, font=\footnotesize, anchor=north west] at ($(p)+(0.06,-0.1)$) {$+q$};
  \node[etiqf, font=\footnotesize, anchor=south east] at ($(m)+(-0.06,0.1)$) {$-q$};
  \node[etiq] at (0,-2.15) {$\vtau=\vp\times\vE$\quad($\otimes$, entrante)};
\end{tikzpicture}
\caption{Par de fuerzas sobre $\vp$}
\end{subfigure}\hfill
\begin{subfigure}[b]{0.29\textwidth}\centering
\begin{tikzpicture}[scale=0.82]
  \fill[dielec, rounded corners=2pt] (0,0) rectangle (3.2,3.2);
  \pgfmathsetseed{17}
  \foreach \i in {0,...,5} \foreach \j in {0,...,5} {
    \pgfmathsetmacro{\ang}{360*rnd}
    \draw[campoP, line width=0.75pt, -{Stealth[length=1.5mm,width=1.2mm]}]
      ({0.32+0.51*\i-0.17*cos(\ang)},{0.32+0.51*\j-0.17*sin(\ang)}) --
      ({0.32+0.51*\i+0.17*cos(\ang)},{0.32+0.51*\j+0.17*sin(\ang)});
  }
  \node[etiq] at (1.6,3.55) {$\vE=\vcero$};
\end{tikzpicture}
\caption{$\langle\vp\rangle=\vcero$}
\end{subfigure}\hfill
\begin{subfigure}[b]{0.29\textwidth}\centering
\begin{tikzpicture}[scale=0.82]
  \fill[dielec, rounded corners=2pt] (0,0) rectangle (3.2,3.2);
  \pgfmathsetseed{17}
  \foreach \i in {0,...,5} \foreach \j in {0,...,5} {
    \pgfmathsetmacro{\ang}{110*(rnd-0.5)}
    \draw[campoP, line width=0.75pt, -{Stealth[length=1.5mm,width=1.2mm]}]
      ({0.32+0.51*\i-0.17*cos(\ang)},{0.32+0.51*\j-0.17*sin(\ang)}) --
      ({0.32+0.51*\i+0.17*cos(\ang)},{0.32+0.51*\j+0.17*sin(\ang)});
  }
  \draw[vecE] (0.7,3.55) -- (2.5,3.55) node[right, etiq] {$\vE$};
\end{tikzpicture}
\caption{$\langle\vp\rangle\neq\vcero$}
\end{subfigure}
\caption{Sustancias polares. (a) En un campo uniforme la fuerza neta sobre un dipolo es nula, pero
el par $\vtau=\vp\times\vE$ tiende a alinearlo con el campo. (b) Sin campo, la agitación térmica
orienta los dipolos al azar y el momento medio es nulo. (c) Con campo, la orientación es parcial:
el par compite con la agitación térmica, y el momento medio, no nulo, tiene la dirección de $\vE$.}
\label{fig:orientacion}
\end{figure}


\begin{observacion}[Polarizabilidad de orientación]\label{obs:langevin}
La mecánica estadística (teoría de Langevin y Debye) da, para $pE\ll k_BT$, un momento medio
$\langle\vp\rangle=\dfrac{p^2}{3k_BT}\,\vE$~\cite{Purcell,Jackson}. Para el agua a temperatura ambiente,
$p^2/3k_BT\simeq\SI{3e-39}{F\,m^2}$, unas cuarenta veces la polarizabilidad electrónica del hidrógeno: por
eso las sustancias polares responden mucho más intensamente al campo, y por eso la permitividad relativa
del agua es tan grande ($\er\simeq80$; \cref{tab:permitividades}).
\end{observacion}

\subsection{La aproximación dipolar}\label{sec:aprox-dipolar}

Tanto si la sustancia es polar como si no, el resultado es el mismo: un dieléctrico en un campo eléctrico
adquiere, en cada elemento de volumen macroscópico, un momento dipolar neto en general no nulo.

\begin{definicion}[Aproximación dipolar]\label{def:aprox-dipolar}
En la \emph{aproximación dipolar}, un dieléctrico polarizado se describe como una distribución continua de
dipolos puntuales: cada elemento de volumen $\dd V'$ centrado en $\vr\,'$, pequeño a escala macroscópica pero
con muchas moléculas, posee un momento dipolar $\dd\vp$, y su contribución al potencial en $\vr$ es la de un
dipolo puntual~\cite{T2},
\begin{equation}\label{eq:dphi-dipolo}
  \dd\phi(\vr)=\frac{1}{4\pi\eo}\frac{\dd\vp\cdot\uv R}{R^2}.
\end{equation}
Se desprecian los multipolos de orden superior de cada elemento, cuya contribución decae más deprisa con la
distancia y es irrelevante a escala macroscópica.
\end{definicion}
